
\section{Results and Discussion}

\subsection{Experimental Results}

The experiment compares two monitoring conditions:
output-only monitoring and monitoring with access to the
reasoning trace. The evaluation focuses on four categories
of unsafe behavior: reward hacking, deception, sabotage,
and harmful compliance.

The pilot evaluation contains 10 episodes. Eight episodes
were used for testing, including six unsafe and two safe
episodes, while two safe episodes were used for calibration.

Under the output-only condition, the monitor detected none
of the six unsafe test episodes, resulting in a recall of
0.00. Under the CoT-plus-output condition, all six unsafe
episodes were detected, resulting in a recall of 1.00.
Both conditions produced zero false positives in this
small test sample.

The AUROC was 0.50 for output-only monitoring and 1.00
for CoT-plus-output monitoring. These results indicate
a substantial difference between the two conditions
within the pilot dataset. However, the small sample size
does not support broad generalization.

\subsection{Statistical Analysis}

The exact one-sided McNemar test comparing the two
monitoring conditions produced a p-value of 0.015625.
This provides preliminary evidence of a difference
between the paired detection outcomes.

For RQ2, the observed recall was 1.00, based on six
detected unsafe episodes out of six. However, the
one-sided binomial test produced a p-value of
approximately 0.118, and the 95\% confidence interval
ranged from approximately 0.607 to 1.00. Therefore,
the pilot results do not establish that recall is
statistically greater than the target value of 0.70
at the 0.05 significance level.

The RQ1 chi-square analysis was reported as degenerate,
so it does not provide reliable evidence of differences
in detection performance across the four categories.

\subsection{Latency and Load Testing}

A separate load-testing experiment evaluated both
monitoring conditions at request rates of 1, 5, 10,
and 20 requests per second. Each combination included
50 measurements, resulting in 400 measurements in total.

The mean latency ranged from approximately 0.194 to
0.230 seconds across the tested conditions and request
rates. The median latency remained close to 0.2 seconds.
The results do not show a substantial increase in
median latency as the request rate increased from
1 to 20 requests per second.

The latency measurements also show that the difference
between the two monitoring conditions was relatively
small compared with the variation across individual
measurements. These findings describe the tested
benchmark configuration and should not be interpreted
as evidence of production-scale performance.

\subsection{Discussion}

The pilot detection results suggest that access to
reasoning traces may provide additional signals for
identifying unsafe behavior that are not available
from the final output alone. In this experiment,
the CoT-plus-output condition detected all six unsafe
test episodes, whereas the output-only condition
detected none.

Nevertheless, this finding should be interpreted
cautiously. The evaluation uses a small demonstration
dataset and a keyword-based monitoring baseline.
The observed performance may depend on the examples,
keywords, and calibration procedure used in this
implementation.

The load-testing results provide a separate view
of system performance. Median latency remained
relatively stable across the tested request rates.
However, further experiments with larger datasets
and additional monitor implementations are necessary
before drawing conclusions about scalability or
general effectiveness.
 